\chapter{Group A Streptococcus}

\section*{What Is This Test?}
Group A Streptococcus (GAS) testing detects \textit{Streptococcus pyogenes}, a gram-positive, catalase-negative, beta-hemolytic coccus that forms chains and is classified by Lancefield group A carbohydrate antigen. GAS causes a spectrum of disease ranging from mild to life-threatening: acute pharyngitis ("strep throat"), scarlet fever (pharyngitis with erythrogenic toxin-mediated diffuse erythematous "sandpaper" rash), impetigo (superficial pyoderma), erysipelas/cellulitis, necrotizing fasciitis (type II, synergistic with anaerobes), streptococcal toxic shock syndrome (STSS, hypotension + multiorgan failure), puerperal sepsis, and bacteremia. Nonsuppurative (immune-mediated) sequelae include acute rheumatic fever (ARF, with carditis, polyarthritis, chorea, erythema marginatum, subcutaneous nodules) and post-streptococcal glomerulonephritis (PSGN). Prompt treatment of GAS pharyngitis (within 9 days of symptom onset) prevents ARF, but does not prevent PSGN. Testing for acute GAS pharyngitis focuses on rapid antigen detection tests (RADTs) and throat culture. Antibody testing (anti-streptolysin O/ASO, anti-DNase B) is used for retrospective diagnosis of recent GAS infection when complications (ARF, PSGN) are suspected weeks after the acute illness.

\section*{Alternative Names}
Strep Test; Rapid Strep Test; RADT; Strep Antigen; Throat Culture; GAS Culture; ASO Titer; Anti-DNase B; Streptozyme

\section*{Principle}
Rapid antigen detection test (RADT): Lateral-flow immunochromatographic assay detecting GAS Lancefield group A carbohydrate antigen from throat swab specimens. Enzyme-based extraction releases the antigen from the bacterial cell wall. Sensitivity 70--90\%; specificity 95--99\%. Results available in 5--15 minutes. Throat culture: Inoculation of a throat swab onto 5\% sheep blood agar with a trimethoprim-sulfamethoxazole (SXT) or bacitracin (Taxo A) disk. Incubate at 35--37\textdegree{}C for 18--48 hours under 5--10\% CO$_2$ or anaerobically. Beta-hemolytic colonies (clear zone of hemolysis) that are catalase-negative, gram-positive cocci in chains, and susceptible to bacitracin (zone of inhibition $\geq$15 mm with 0.04-unit disk) are identified as \textit{S. pyogenes}. Confirmatory identification by Lancefield grouping (latex agglutination) or MALDI-TOF. Throat culture is 90--95\% sensitive (reference standard). ASO titer: Neutralization assay or nephelometric/turbidimetric immunoassay measuring antibodies against streptolysin O, a hemolytic exotoxin. ASO rises 1--3 weeks after acute GAS infection, peaks at 3--6 weeks, and declines over months. A single elevated titer (>200 IU/mL in adults; >300 IU/mL in children) or a fourfold rise between acute and convalescent sera (2--4 weeks apart) indicates recent GAS infection. Upper limit of normal varies by assay and laboratory. Anti-DNase B: Antibody against deoxyribonuclease B, another GAS exotoxin. Rises more slowly than ASO and peaks later. Anti-DNase B is useful when ASO fails to rise (which occurs in 20--30\% of patients with streptococcal pharyngitis, especially in skin infections---streptolysin O is inhibited by skin lipids, so ASO often does not rise after GAS impetigo).

\section*{History}
\textit{Streptococcus pyogenes} was first described by Theodor Billroth in 1874 and named \textit{Streptococcus} (twisted berry) by Friedrich Fehleisen in 1883. The Lancefield grouping system was developed by Rebecca Lancefield in the 1930s. The association between GAS pharyngitis and rheumatic fever was established in the 1930s--1940s. Penicillin was shown to prevent ARF by Floyd Denny and Lewis Wannamaker in the 1950s. Rapid antigen detection tests were introduced in the 1980s, and their use in combination with a reflex throat culture for negative RADTs became the standard of care. Clinician decision rules (Centor criteria, modified for age: McIsaac score) were developed to stratify the pretest probability of GAS pharyngitis and guide testing. The IDSA updated its guideline for diagnosis and management of GAS pharyngitis in 2012, recommending RADT as the initial diagnostic approach with reflex throat culture for negative RADTs in children and adolescents (but not necessarily in adults with low Centor scores).

\section*{Why This Test Is Ordered}
\begin{itemize}
  \item Rapid diagnosis of GAS pharyngitis in patients presenting with acute sore throat and Centor/McIsaac criteria (tonsillar exudates, tender anterior cervical lymphadenopathy, fever >38\textdegree{}C, absence of cough, age 3--14 years earns 1 point; $\geq$3 points warrants RADT/culture)
  \item Differentiation of bacterial (GAS) pharyngitis from viral pharyngitis to guide appropriate antibiotic therapy and reduce unnecessary antibiotic use for viral infections (the majority of pharyngitis cases). Empiric antibiotics based on clinical criteria alone should be avoided except in very specific circumstances
  \item Prevention of acute rheumatic fever: Appropriate antibiotic therapy for GAS pharyngitis within 9 days of symptom onset prevents ARF, which remains a major cause of acquired heart disease in children in resource-limited settings
  \item Prevention of suppurative complications: peritonsillar abscess, retropharyngeal abscess, cervical lymphadenitis, acute otitis media, sinusitis
  \item Investigation of suspected invasive GAS infection: Blood cultures, tissue cultures, and operative specimens from patients with necrotizing fasciitis, STSS, or GAS bacteremia
  \item ASO/anti-DNase B antibody titers for retrospective diagnosis of recent GAS infection in patients with suspected ARF, PSGN, or pediatric autoimmune neuropsychiatric disorders associated with streptococcal infections (PANDAS)
  \item Screening of close contacts of patients with invasive GAS disease (necrotizing fasciitis, STSS), though evidence for chemoprophylaxis in contacts is limited and not routinely recommended
\end{itemize}

\section*{Primary vs Secondary Test}
RADT and throat culture are primary diagnostic tests for acute GAS pharyngitis. A negative RADT in a child or adolescent should be followed by a reflex throat culture (IDSA strong recommendation, high-quality evidence) because of the greater consequence of missing GAS. In adults, especially with low Centor scores, a negative RADT alone is sufficient to rule out GAS without culture due to the lower risk of ARF and suppurative complications. ASO/anti-DNase B titers are secondary (retrospective) tests used when the acute illness has resolved.

\section*{How This Test Is Performed}
Throat swab collection is critical for accuracy: Using a tongue depressor, a Dacron or rayon-tipped swab is vigorously rubbed over both tonsillar pillars, the posterior pharynx, and any exudative or erythematous areas. The tongue, uvula, and buccal mucosa should be avoided. The swab should be placed immediately in the RADT kit or transported in Amies transport medium (for culture). Throat culture specimens should be transported to the laboratory within 24 hours at room temperature. For ASO/anti-DNase B testing, venous blood in an SST tube is standard; acute and convalescent sera should be drawn 2--4 weeks apart for paired titers.

\section*{What Preparation Is Needed}
No special preparation for throat swab. Patients should avoid eating, drinking, or using mouthwash for at least 30 minutes before throat swab collection. Recent antibiotic use reduces RADT and throat culture sensitivity; if possible, testing should be performed before antibiotics are initiated. Antibiotic therapy should never be delayed in a critically ill patient (suspected necrotizing fasciitis, STSS) while awaiting test results.

\section*{Contraindications and Precautions}
RADT sensitivity is operator-dependent and may be reduced with poorly collected specimens (inadequate swabbing of the tonsillar pillars and posterior pharynx). In children and adolescents, a negative RADT should be confirmed by throat culture because the sensitivity of RADT is 70--90\% and because children have a higher prevalence of GAS pharyngitis and are at greater risk for ARF. GAS pharyngeal carriage (asymptomatic colonization, 5--20\% of school-age children during winter/spring) will test positive by RADT and culture. When a child with GAS carriage develops a viral pharyngitis, the RADT/culture will be positive but the clinical illness is due to a virus, not GAS. Testing should be targeted to patients with clinical signs and symptoms of GAS pharyngitis (Centor score $\geq$2--3). Testing asymptomatic individuals or those with predominantly viral symptoms (cough, rhinorrhea, conjunctivitis, hoarseness, diarrhea) should be avoided, as a positive result likely reflects carriage and may lead to unnecessary antibiotic exposure. ASO titer interpretation: A single elevated ASO titer is not diagnostic of ARF; the Jones criteria require evidence of preceding GAS infection (positive throat culture, elevated or rising ASO or anti-DNase B) plus clinical major/minor criteria to diagnose ARF. ASO titers have a wide normal range that varies by age, season, and geography. Elevated ASO may persist for months after acute GAS infection and does not distinguish recent from remote GAS infection.

\section*{Risks}
Throat swabs cause transient gagging and discomfort but have no significant risks. Venipuncture for ASO testing has standard risks. The major risk is failure to diagnose and treat GAS pharyngitis in a child who subsequently develops ARF with permanent cardiac valve damage.

\section*{What Happens After the Test}
RADT results in 5--15 minutes. Throat culture preliminary at 18--24 hours; final negative at 48 hours. If RADT is positive, the patient should be treated with penicillin V (250 mg BID-TID for children, 500 mg BID for adults) or amoxicillin (50 mg/kg once daily for children, max 1,000 mg; 1,000 mg once daily for adults) for 10 days. First-generation cephalosporins, clindamycin, or macrolides are alternatives for penicillin-allergic patients. Treatment reduces symptom duration by 12--24 hours if started early (within 48 hours of onset), reduces risk of ARF by >80\%, reduces suppurative complications, and reduces infectivity (patients are considered noninfectious after 24 hours of antibiotics). Patients should be counseled to complete the full 10-day course despite symptom improvement.

\section*{Understanding Results}

\subsection*{Below Normal}
\begin{itemize}
  \item \textbf{RADT negative + throat culture negative:} No GAS detected. The pharyngitis is presumed viral. Symptomatic treatment (analgesia, hydration, throat lozenges) is appropriate; antibiotics are not indicated. If symptoms persist beyond 5--7 days or worsen, consideration of alternative diagnoses (peritonsillar abscess, epiglottitis, Lemierre syndrome, HIV seroconversion, EBV, CMV) is warranted
  \item \textbf{ASO titer normal (<200 IU/mL adult; <300 IU/mL child):} No evidence of recent GAS infection. A normal ASO does not entirely exclude recent GAS infection (ASO fails to rise in 20--30\% of GAS pharyngitis, particularly after skin infections). Anti-DNase B levels should be checked if clinical suspicion remains high
\end{itemize}

\subsection*{Above Normal}
\begin{itemize}
  \item \textbf{RADT positive:} GAS antigen detected. In a symptomatic patient, this is diagnostic of GAS pharyngitis. Treat with penicillin or amoxicillin. No confirmatory culture is needed. In an asymptomatic patient or one with predominantly viral symptoms (cough, rhinorrhea), the result likely reflects GAS carriage rather than acute infection, and antibiotics are generally not recommended
  \item \textbf{Throat culture positive for \textit{Streptococcus pyogenes}:} Definitive diagnosis. Same interpretation as RADT positive; correlation with clinical presentation is essential to distinguish infection from carriage
  \item \textbf{ASO titer elevated (>200 IU/mL adult; >300 IU/mL child) or fourfold rise between acute and convalescent sera:} Evidence of recent GAS infection. In a patient with suspected ARF (Jones criteria), this supports the diagnosis (evidence of preceding GAS infection). In a patient with suspected PSGN, an elevated ASO or anti-DNase B combined with acute nephritic syndrome (hematuria, proteinuria, edema, hypertension, low C3) supports PSGN. ASO titers do not guide treatment of acute pharyngitis (they rise too late); they are retrospective diagnostic tools
  \item \textbf{Anti-DNase B elevated:} Evidence of recent GAS infection. The combination of ASO + anti-DNase B increases diagnostic sensitivity to >90\% for both pharyngeal and skin GAS infections. If both are negative, ARF is unlikely (though rare cases occur)
\end{itemize}

\section*{Normal Results}
RADT: \textbf{Negative}. Throat culture: \textbf{No Group A Streptococcus (Streptococcus pyogenes) isolated} or \textbf{Normal respiratory flora}. ASO titer: \textbf{<200 IU/mL} (adult) or \textbf{<300 IU/mL} (child) --- varies by assay. Anti-DNase B: \textbf{Within normal limits} (varies by laboratory).

\section*{Follow-up Tests}
\begin{itemize}
  \item \textbf{Throat culture:} For children and adolescents with negative RADT, reflex to throat culture (IDSA recommendation). Culture should be followed and, if positive, the patient should be called and antibiotics initiated
  \item \textbf{ASO and anti-DNase B titers (paired, acute and convalescent):} For suspected ARF, PSGN, or PANDAS. A fourfold rise between acute (first 1--2 weeks) and convalescent (4--6 weeks) serum is highly specific for recent GAS
  \item \textbf{Blood cultures (2 sets):} For suspected invasive GAS (necrotizing fasciitis, STSS, bacteremia). Blood cultures are positive in 60--70\% of STSS cases
  \item \textbf{Echocardiogram:} For suspected ARF, to evaluate for carditis (mitral and/or aortic regurgitation, pericardial effusion, valvular dysfunction) per the Jones criteria
  \item \textbf{Complete blood count, CRP/ESR:} Inflammatory markers are elevated in ARF and support the diagnosis (minor Jones criteria)
  \item \textbf{Urinalysis and complement C3/C4:} For suspected PSGN: hematuria (often macroscopic), proteinuria, RBC casts, and low C3 (returns to normal within 6--8 weeks in PSGN; persistent hypocomplementemia suggests another cause of glomerulonephritis)
\end{itemize}

\section*{References}
\begin{enumerate}
  \item Shulman ST, Bisno AL, Clegg HW, et al. Clinical practice guideline for the diagnosis and management of group A streptococcal pharyngitis: 2012 update by the IDSA. Clin Infect Dis. 2012;55(10):e86-e102.
  \item Gewitz MH, Baltimore RS, Tani LY, et al. Revision of the Jones criteria for the diagnosis of acute rheumatic fever in the era of Doppler echocardiography. Circulation. 2015;131(20):1806-1818.
  \item Carapetis JR, Steer AC, Mulholland EK, Weber M. The global burden of group A streptococcal diseases. Lancet Infect Dis. 2005;5(11):685-694.
  \item Stevens DL, Bryant AE. Necrotizing soft-tissue infections. N Engl J Med. 2017;377(23):2253-2265.
  \item Kanwal S, Vaitla P, Aung AC, et al. Strep pharyngitis testing. StatPearls; 2023.
\end{enumerate}

\clearpage
